\section{Robustness to the Taylor-Rule Reaction Coefficients}
\label{sec:robustness}

The quantitative magnitude of the impulse responses analysed in
Section~\ref{sec:results} hinges on the central-bank reaction
coefficients $\phi_{\pi}$ and $\phi_{y}$. To assess the sensitivity of
the baseline results, I recompute the policy function and the impulse
responses under three alternative configurations: a hawkish rule
($\phi_{\pi}=2.0$ with $\phi_{y}=0.125$), a dovish rule
($\phi_{\pi}=1.1$ with $\phi_{y}=0.125$), and a strict
inflation-targeting rule ($\phi_{\pi}=1.5$ with $\phi_{y}=0$).

Two observations follow. First, the qualitative pattern of the impulse
responses is preserved across all four calibrations: output and
inflation fall together, the nominal rate rises, and the real rate
rises by more than the nominal rate. The Blanchard--Kahn condition
continues to hold in every case because $\phi_{\pi}>1$ ensures
determinacy through the Taylor principle.

Second, the quantitative magnitudes vary in economically intuitive ways.
A more hawkish rule (higher $\phi_{\pi}$) attenuates the on-impact
decline in output because the larger nominal-rate response makes the
real-rate increase shorter-lived in expectation; the inflation response
correspondingly becomes more muted. A dovish rule has the opposite
effect: a smaller $\phi_{\pi}$ implies a smaller nominal-rate response,
a longer expected period of below-target inflation, and therefore
larger declines in output and inflation. The strict
inflation-targeting variant ($\phi_{y}=0$) leaves the impulse responses
quantitatively close to the baseline because $\phi_{y}=0.125$ in the
baseline is already small.

Together, these exercises confirm that the qualitative properties of
the baseline impulse response are not artefacts of the specific values
of the Taylor-rule reaction coefficients, while also highlighting that
the quantitative magnitudes vary with the calibrated central-bank
preferences in a manner consistent with the analytical predictions of
the model.
